\documentclass[12pt]{article}

%% ---------------------------------------------------------------
%% Packages
%% ---------------------------------------------------------------
\usepackage[margin=1in]{geometry}
\usepackage{amsmath, amssymb, amsthm}
\usepackage{booktabs}
\usepackage{hyperref}
\usepackage{graphicx}
\usepackage{xcolor}
\usepackage{microtype}
\usepackage{enumitem}
\usepackage{algorithm}
\usepackage{algpseudocode}
\usepackage{natbib}

%% ---------------------------------------------------------------
%% Theorem environments
%% ---------------------------------------------------------------
\newtheorem{definition}{Definition}
\newtheorem{remark}{Remark}
\newtheorem{observation}{Observation}

%% ---------------------------------------------------------------
%% Custom commands
%% ---------------------------------------------------------------
\newcommand{\Qm}{Q_m}
\newcommand{\Pm}{P_m}
\newcommand{\dH}{d_H}
\newcommand{\Bm}{B(m,r)}
\newcommand{\E}{\mathbb{E}}
\newcommand{\Z}{\mathbb{Z}}
\newcommand{\N}{\mathbb{N}}

%% ---------------------------------------------------------------
%% Metadata
%% ---------------------------------------------------------------
\title{%
  Prime Search in the Odd Binary Hypercube:\\
  Number-Line Locality, Hamming-Space Locality, and Pareto Tradeoffs
}

\author{
  Preetam Ojha\\
  \texttt{nehu.preet@gmail.com}
}

\date{May 2026}

%% arXiv subject classification
%% Primary:   cs.DM  (Discrete Mathematics)
%% Secondary: cs.DS  (Data Structures and Algorithms)
%% Suggested cross-list: math.CO, math.NT

%% ---------------------------------------------------------------
\begin{document}
%% ---------------------------------------------------------------

\maketitle

\begin{center}
\small
\textbf{Suggested arXiv subject classes:} cs.DM (primary), cs.DS (secondary);
possible cross-list: math.CO, math.NT
\end{center}

\begin{abstract}
We study prime numbers as vertices in the $m$-bit binary hypercube
$Q_m = \{0,1\}^m$, where adjacency is defined by Hamming distance rather than
arithmetic distance. This yields a search geometry fundamentally different from
the classical number line. We ask: given an integer $x$, how close is it to a
prime in Hamming distance, and how does Hamming-space search compare with
standard odd-increment prime search?

We report four empirical findings. \emph{First}, experiments show that nearly
every $m$-bit integer is within Hamming distance 2 of a prime for the tested
small dimensions, consistent with the rapid growth of Hamming balls relative
to the decay of prime density. \emph{Second}, in the full hypercube, bit~0
(parity) overwhelmingly dominates nearest-prime bit-flip paths. This motivates
restricting attention to the \emph{odd subcube}, where bit~0 is fixed and only
odd candidates are searched. \emph{Third}, after this parity correction,
standard odd-increment search achieves the lowest average candidate-check count
among the tested strategies, but odd-subcube Hamming search wins or ties
standard search on approximately 58--60\% of individual starting points across
$m \in \{24,28,32\}$. \emph{Fourth}, Pareto analysis shows that the two search
families optimize different objectives: standard search finds arithmetically
close primes, while Hamming search finds bit-close primes.

Thus, the main contribution is not a replacement for standard prime search, but
an experimental comparison of two geometries of primality: number-line locality
and Hamming-space locality.

We release all code and data at
\url{https://github.com/pojhafb/prime-hypercube-search}.
\end{abstract}

\tableofcontents
\newpage

%% ---------------------------------------------------------------
\section{Introduction}
\label{sec:intro}
%% ---------------------------------------------------------------

Prime numbers are classically studied as points on the integer number line,
where distance is measured by subtraction: $|p-x|$. The prime number theorem
\citep{hadamard1896, delavalleepoussin1896} tells us that primes are sparse:
among integers near $N$, about $1/\ln N$ of them are prime. The expected gap to
the next prime from a typical integer near $x$ is approximately $\ln x$.

Every integer also has a binary representation, and binary strings of length
$m$ are naturally placed as vertices of the $m$-dimensional hypercube
$Q_m = \{0,1\}^m$. Two vertices are adjacent if they differ in exactly one bit.
The number of differing bit positions between two integers $x$ and $y$ is the
\emph{Hamming distance}:
\[
  \dH(x,y) = |\{i : \mathrm{bit}_i(x) \neq \mathrm{bit}_i(y)\}|.
\]
This defines a metric on $Q_m$ that is very different from arithmetic distance.

This paper asks: \emph{what does the distribution of primes look like in
Hamming geometry?} More concretely, can one find a prime near a given integer
$x$ by searching within a small Hamming ball around $x$, and how does such
search compare to the standard algorithm of testing $x, x+2, x+4, \ldots$?

\subsection{Why this question is interesting}

There are two independent motivations.

\paragraph{Geometric.}
The primes are a fixed, arithmetically-defined set. Placing them in two
different metric spaces---the integers $\Z$ with ordinary absolute value, and
the hypercube $Q_m$ with Hamming distance---and comparing their local density
structure is a natural question in discrete mathematics. It draws on the
combinatorics of the hypercube \citep{macwilliams1977, knuth2011} and is
loosely motivated by the study of prime gaps \citep{granville1995}, though the
connection here is geometric rather than analytic.

\paragraph{Computational.}
Finding a prime near a given large integer is a basic primitive in
cryptographic key generation \citep{menezes1996}. Standard algorithms test
$x, x+2, x+4, \ldots$ and rely on the prime number theorem for their expected
cost. An alternative search in Hamming space produces primes with different
characteristics---in particular, primes that may be arithmetically far from
$x$ but differ in only a few bits. Whether bit-close primes have practical
advantages over arithmetic-close primes is an open question; the goal here is
to characterize the two search geometries rather than to advocate for one.

\subsection{Summary of results}

\begin{enumerate}[label=(\roman*)]
  \item \textbf{Hamming density (empirical).} In a full scan of $Q_{14}$ and
    sampled experiments at $m \in \{16,18,20\}$, nearly every tested integer is
    within a small Hamming radius of a prime. In the $m=14$ scan, 98.4\% of
    integers are within distance 2 and 100\% within distance 3.

  \item \textbf{Parity dominates the full hypercube.} In unrestricted Hamming
    search, bit~0 dominates the flip-position distribution because flipping it
    changes parity. Since all primes greater than 2 are odd, this motivates
    the odd-subcube formulation.

  \item \textbf{Learned policies recover low-order structure.} A bit-order
    policy learned from data transfers across dimensions. In the full hypercube,
    much of the learned signal comes from parity. After restricting to the odd
    subcube, learned policies behave similarly to low-bit-first ordering,
    suggesting that the strongest effect is due to parity and low-order modular
    structure.

  \item \textbf{SOI wins average checks; LBF is competitive per-instance.}
    In the odd subcube, standard odd-increment search achieves the lowest
    average candidate-check count across all tested dimensions. However,
    low-bit-first Hamming search wins or ties standard search on roughly
    58--60\% of individual starting points.

  \item \textbf{Two geometries of prime search.} Standard search finds
    arithmetically close primes; Hamming search finds bit-close primes. Pareto
    analysis shows that the strategies optimize different objectives rather
    than one simply dominating the other.
\end{enumerate}

\subsection{Organization}

Section~\ref{sec:model} sets up the mathematical model.
Section~\ref{sec:strategies} describes the search strategies.
Section~\ref{sec:theory} provides heuristic estimates.
Section~\ref{sec:experiments} presents experimental results.
Section~\ref{sec:pareto} gives the Pareto analysis.
Section~\ref{sec:related} discusses related work.
Section~\ref{sec:conclusion} concludes with limitations and open questions.

%% ---------------------------------------------------------------
\section{Mathematical Model}
\label{sec:model}
%% ---------------------------------------------------------------

\subsection{The binary hypercube}

\begin{definition}
  The \emph{$m$-bit binary hypercube} is the graph
  $Q_m=(\{0,1\}^m,E)$ where $\{x,y\}\in E$ if and only if
  $\dH(x,y)=1$.
\end{definition}

We identify each vertex with the integer it encodes in binary, so $Q_m$
contains the integers $\{0,1,\ldots,2^m-1\}$.

\begin{definition}
  The \emph{Hamming ball} of radius $r$ centered at $x$ is:
  \[
    \mathcal{B}(x,r)=\{y\in Q_m : \dH(x,y)\leq r\}.
  \]
  Its size depends only on $m$ and $r$:
  \[
    |\mathcal{B}(x,r)| = B(m,r)=\sum_{i=0}^{r}\binom{m}{i}.
  \]
\end{definition}

\subsection{Primes in the hypercube}

\begin{definition}
  The \emph{prime set} of $Q_m$ is
  \[
    P_m=\{x\in Q_m : x \text{ is prime}\}.
  \]
  The \emph{Hamming distance to the nearest prime} from $x$ is:
  \[
    D_m(x)=\min_{p\in P_m}\dH(x,p).
  \]
\end{definition}

By the prime number theorem, $|P_m| \approx 2^m/(m\ln 2)$, so the prime density
in $Q_m$ is approximately:
\[
  \delta_m \approx \frac{1}{m\ln 2}.
\]

\subsection{The odd subcube}

All primes greater than 2 are odd. Therefore, for $m\geq 2$, all nontrivial
prime vertices lie in:
\[
  Q_m^{\mathrm{odd}}
  =
  \{x\in Q_m : \mathrm{bit}_0(x)=1\}.
\]
$Q_m^{\mathrm{odd}}$ is the odd subcube, a sub-hypercube of dimension $m-1$
with $2^{m-1}$ vertices. In this paper, odd-subcube Hamming search fixes bit~0
to 1 and flips only bits $1,\ldots,m-1$.

%% ---------------------------------------------------------------
\section{Search Strategies}
\label{sec:strategies}
%% ---------------------------------------------------------------

Given a starting integer $x$, we set:
\[
  x_0 = x\,|\,1,
\]
that is, $x_0$ is the odd integer obtained by setting bit~0 to 1. A search
strategy returns a prime $p$, a candidate-check count $k$, the Hamming distance
$\dH(x_0,p)$, and the arithmetic distance $|p-x_0|$.

We compare six strategies.

\subsection{Standard odd-increment search (SOI)}

\begin{algorithm}[H]
\caption{Standard odd-increment search}
\begin{algorithmic}[1]
\State $x_0 \gets x \,|\, 1$
\State $k \gets 0$, $c \gets x_0$
\Loop
  \State $k \gets k + 1$
  \If{$c$ is prime}
    \Return $(c,\, k,\, \dH(x_0,c),\, c-x_0)$
  \EndIf
  \State $c \gets c+2$
\EndLoop
\end{algorithmic}
\end{algorithm}

This is the standard arithmetic-local algorithm. It always finds the smallest
prime $\geq x_0$.

\subsection{Odd-subcube Hamming search}

Let $\mathcal{F}(m,r_{\max})$ be a list of flip sets---tuples of bit positions
in $\{1,\ldots,m-1\}$---ordered by some policy $\pi$.

\begin{algorithm}[H]
\caption{Odd-subcube Hamming search}
\begin{algorithmic}[1]
\State $x_0 \gets x \,|\, 1$
\State $k \gets 0$
\For{each flip set $S \in \mathcal{F}(m,r_{\max})$ ordered by $\pi$}
  \State $k \gets k+1$
  \State $c \gets x_0 \oplus \bigoplus_{b\in S}2^b$
  \If{$c$ is prime}
    \Return $(c,\, k,\, |S|,\, |c-x_0|)$
  \EndIf
\EndFor
\end{algorithmic}
\end{algorithm}

The choice of ordering $\pi$ gives the following concrete strategies:

\begin{itemize}
  \item \textbf{LBF} (\emph{low-bit-first}): sort flip sets by
    $(\sum S,\max S)$, preferring lower bit indices.
  \item \textbf{HBF} (\emph{high-bit-first}): sort by
    $(-\sum S,-\max S)$.
  \item \textbf{RND} (\emph{uniform random}): shuffle flip sets randomly.
  \item \textbf{LRN} (\emph{learned bit-order}): sort by learned bit scores.
\end{itemize}

\subsection{Hybrid strategy}

The hybrid strategy interleaves SOI candidates and learned Hamming candidates,
deduplicating before testing.

\begin{algorithm}[H]
\caption{Hybrid standard + learned search}
\begin{algorithmic}[1]
\State $x_0 \gets x \,|\, 1$, $\mathrm{seen}\gets \emptyset$, $k\gets 0$
\For{$i=0,1,2,\ldots$}
  \If{$i<L$} add $x_0+2i$ to candidates \EndIf
  \If{$i<|\mathcal{F}|$} add $x_0\oplus \mathrm{mask}(\mathcal{F}_i)$ to candidates \EndIf
  \For{each candidate $c$ not in $\mathrm{seen}$}
    \State $\mathrm{seen}\gets \mathrm{seen}\cup\{c\}$, $k\gets k+1$
    \If{$c$ is prime}
      \Return $(c,k,\dH(x_0,c),|c-x_0|)$
    \EndIf
  \EndFor
\EndFor
\end{algorithmic}
\end{algorithm}

$L$ is the odd-increment limit, set to 128 in the reported experiments.

\subsection{Learned bit-order policy}
\label{sec:learned}

We learn per-bit usefulness scores from data. Given training dimension
$m_{\mathrm{tr}}$ and $N_{\mathrm{tr}}$ random odd starting points:

\begin{enumerate}
  \item Run LBF on each starting point and record the successful flip set $S^*$.
  \item For each bit $b\in S^*$, increment $\mathrm{score}[b]$.
  \item Bits with zero count receive a fallback score equal to
    $0.25$ times the minimum positive score.
  \item Bit~0 is forced to score 0 and is excluded from odd-subcube search.
\end{enumerate}

For test dimension $m_{\mathrm{te}}>m_{\mathrm{tr}}$, learned scores are
transferred by retaining the scores for old bit positions and assigning the
fallback score to new higher bit positions.

%% ---------------------------------------------------------------
\section{Heuristic Analysis}
\label{sec:theory}
%% ---------------------------------------------------------------

This section provides heuristic estimates and informal arguments. No new
theorems are proved.

\subsection{Expected checks for SOI}

The prime number theorem implies that the average gap between consecutive
primes near $N$ is approximately $\ln N$. Restricting to odd candidates halves
this, giving:
\[
  \E[\text{checks for SOI}]
  \approx
  \frac{\ln(2^m)}{2}
  =
  \frac{m\ln 2}{2}.
\]
For $m=32$, this gives:
\[
  \frac{32\ln 2}{2}\approx 11.1.
\]

\subsection{Hamming ball size and prime density}

The expected number of primes within Hamming radius $r$ of a random odd
$m$-bit integer is approximately:
\[
  \E[|P_m \cap \mathcal{B}(x_0,r)|]
  \approx
  B(m-1,r)\cdot \frac{2}{m\ln 2},
\]
where the factor of 2 accounts for restricting to odd candidates.

For $m=32$ and $r=4$:
\[
  B(31,4)
  =
  1+31+465+4495+31465
  =
  36457,
\]
and:
\[
  \frac{2}{32\ln 2}\approx 0.090.
\]
So the radius-4 odd-subcube ball contains many primes in expectation. The
difficulty is not whether a prime exists nearby in Hamming distance, but how
many ordered candidate checks are needed to reach one.

\subsection{Why SOI dominates average checks}

SOI tests candidates in increasing arithmetic order and directly exploits the
local prime density on the number line. Hamming search visits candidates in
bit-flip order, which does not align with arithmetic locality.

\begin{observation}
A Hamming flip of bit $k$ changes $x$ by $\pm 2^k$, which may be a very large
arithmetic step. Hamming search therefore optimizes bit locality rather than
number-line locality.
\end{observation}

\subsection{Why Hamming search wins on many individual inputs}

\begin{observation}
If a prime occurs early in the Hamming ordering of $x_0$, then Hamming search
can beat SOI for that starting point even when SOI has better average behavior.
\end{observation}

This accounts for the empirical fact that LBF wins or ties SOI on a large
fraction of individual inputs, despite losing on average candidate checks.

%% ---------------------------------------------------------------
\section{Experiments}
\label{sec:experiments}
%% ---------------------------------------------------------------

\subsection{Setup}

All primality tests use \texttt{sympy.isprime}. Results are cached with
\texttt{functools.lru\_cache}. The goal is not to benchmark primality-testing
algorithms themselves, but to compare candidate-generation policies. All
strategies in a given experiment use the same primality-testing backend.

Unless otherwise stated, the reported experiments use:

\begin{itemize}
  \item $m_{\mathrm{tr}}=20$
  \item $m_{\mathrm{te}}\in\{24,28,32\}$
  \item $N_{\mathrm{tr}}=5000$ training samples
  \item $N_{\mathrm{te}}=10000$ test samples per dimension
  \item $r_{\max}=4$
  \item seed $=42$
\end{itemize}

\subsection{Hamming distance distribution}

Table~\ref{tab:dist14} shows the full scan for $m=14$.

\begin{table}[h]
\centering
\caption{Hamming distance from each integer in $Q_{14}$ to its nearest prime.
Full scan, $2^{14}=16384$ integers.}
\label{tab:dist14}
\begin{tabular}{rrrr}
\toprule
$D_{14}(x)$ & Count & \% & Cumulative \% \\
\midrule
0 & 1900 & 11.6 & 11.6 \\
1 & 7933 & 48.4 & 60.0 \\
2 & 6286 & 38.4 & 98.4 \\
3 & 265  & 1.6  & 100.0 \\
$\geq 4$ & 0 & 0.0 & 100.0 \\
\bottomrule
\end{tabular}
\end{table}

This full scan shows that every integer in $Q_{14}$ is within Hamming distance
3 of a prime, and 98.4\% are within distance 2.

\subsection{Average candidate checks}

Table~\ref{tab:checks} reports average candidate checks by strategy.

\begin{table}[h]
\centering
\caption{Average candidate checks by strategy, $m \in \{24,28,32\}$,
$N_{\mathrm{te}}=10000$, $r_{\max}=4$.
SOI = standard odd-increment, LBF = low-bit-first odd-subcube search,
LRN = learned bit-order, HBF = high-bit-first, RND = uniform random.
Lower is better.}
\label{tab:checks}
\begin{tabular}{lrrr}
\toprule
Strategy & $m=24$ & $m=28$ & $m=32$ \\
\midrule
SOI            & \textbf{7.49}  & \textbf{8.73}  & \textbf{10.17} \\
Hybrid         & 7.70  & 9.15  & 10.58 \\
HBF            & 8.35  & 9.76  & 11.22 \\
LRN            & 8.37  & 9.96  & 11.56 \\
LBF            & 8.38  & 9.95  & 11.53 \\
RND            & 8.42  & 9.91  & 11.62 \\
\bottomrule
\end{tabular}
\end{table}

SOI has the fewest average checks at all tested dimensions. Although HBF has
competitive average checks in this table, it does so by testing candidates that
are often arithmetically very far from the starting point. Thus, candidate
checks alone are not a sufficient evaluation metric. Section~\ref{sec:pareto}
evaluates the strategies jointly by checks, Hamming distance, and arithmetic
distance.

\subsection{Per-instance wins}

Table~\ref{tab:wins} shows the pairwise win rate of LBF against SOI.

\begin{table}[h]
\centering
\caption{Per-instance win rate of LBF vs. SOI.
For each starting point $x$, we compare the number of candidate checks used by
each strategy.}
\label{tab:wins}
\begin{tabular}{lrrr}
\toprule
$m$ & LBF wins & SOI wins & Ties \\
\midrule
24 & 37.1\% & 36.6\% & 26.3\% \\
28 & 37.7\% & 39.0\% & 23.3\% \\
32 & 39.2\% & 40.1\% & 20.7\% \\
\bottomrule
\end{tabular}
\end{table}

Thus, LBF wins or ties SOI on 63.4\% ($m=24$), 60.3\% ($m=28$), and 59.9\%
($m=32$) of individual starting points, despite having a higher average check
count. This shows that SOI's average advantage does not imply uniform
per-instance dominance.

\subsection{Distance comparison}

Table~\ref{tab:distances} compares the average Hamming and arithmetic distances
to the found prime for $m=32$.

\begin{table}[h]
\centering
\caption{Average Hamming and arithmetic distance to the found prime, $m=32$.
Values are computed from the odd-subcube experiment with $N_{\mathrm{te}}=10000$
and $r_{\max}=4$.}
\label{tab:distances}
\begin{tabular}{lrr}
\toprule
Strategy & Avg $\dH(x_0,p)$ & Avg $|p-x_0|$ \\
\midrule
SOI    & 2.81 & \textbf{18.34} \\
Hybrid & 1.90 & 665{,}588.18 \\
HBF    & 0.98 & 486{,}532{,}672.18 \\
LBF    & \textbf{0.98} & 24{,}771{,}857.20 \\
LRN    & 0.98 & 25{,}453{,}898.40 \\
RND    & 0.98 & 86{,}389{,}309.05 \\
\bottomrule
\end{tabular}
\end{table}

Odd-subcube Hamming strategies find primes at roughly one bit flip from the
starting point on average. SOI, by contrast, finds primes at an average Hamming
distance of 2.81 bits for $m=32$, but with much smaller arithmetic distance.
This highlights the core tradeoff: SOI is arithmetic-local, while Hamming
search is bit-local.

%% ---------------------------------------------------------------
\section{Pareto Analysis}
\label{sec:pareto}
%% ---------------------------------------------------------------

\subsection{Pareto dominance}

\begin{definition}
Strategy $A$ \emph{Pareto-dominates} strategy $B$ on metrics $(f_1,f_2)$,
both to be minimized, if $f_1(A)\leq f_1(B)$ and $f_2(A)\leq f_2(B)$ with at
least one strict inequality. The \emph{Pareto frontier} is the set of
non-dominated strategies.
\end{definition}

\subsection{Checks vs. Hamming distance}

In the plane:
\[
  (\text{average checks},\ \text{average Hamming distance}),
\]
SOI and odd-subcube Hamming methods represent different tradeoffs. SOI has
the lowest average checks but larger Hamming distance. LBF, LRN, and related
Hamming strategies have slightly larger check counts but much smaller Hamming
distance.

This means neither family uniformly dominates when Hamming distance is included
as an objective.

\subsection{Checks vs. arithmetic distance}

In the plane:
\[
  (\text{average checks},\ \log_2(1+\text{average arithmetic distance})),
\]
SOI dominates most Hamming strategies because it is designed to find the next
prime on the number line. This confirms that SOI is the appropriate strategy
when arithmetic locality is the primary objective.

\subsection{Multi-objective score}

We define a composite score:
\[
  \text{score}
  =
  \overline{k}
  +
  \lambda_H\cdot \overline{\dH}
  +
  \lambda_A\cdot \log_2(1+\overline{|p-x_0|}),
\]
where $\overline{\cdot}$ denotes the average over all starting points.

The best strategy depends on the weights. When $\lambda_A$ is large or
candidate checks are weighted heavily, SOI is favored. When $\lambda_H$ is
large and arithmetic distance is weakly penalized, odd-subcube Hamming
strategies become competitive or preferred. This confirms that the strategies
are best understood as optimizing different geometries rather than as direct
replacements for one another.

%% ---------------------------------------------------------------
\section{Related Work}
\label{sec:related}
%% ---------------------------------------------------------------

\paragraph{Prime gaps.}
The distribution of gaps between consecutive primes is well-studied
\citep{cramer1936, granville1995, maynard2015}. Our work is complementary:
we study not arithmetic gaps but Hamming distances, which encode a different
notion of proximity.

\paragraph{The binary hypercube.}
The hypercube is a classical object in combinatorics and coding theory
\citep{macwilliams1977, knuth2011}. Boolean functions on the hypercube have
been extensively analyzed \citep{odonnell2014}. The distribution of specific
number-theoretic sets, such as primes, as subsets of the hypercube appears less
standard as a unified search-geometry framework.

\paragraph{Primality testing.}
Miller--Rabin \citep{miller1976, rabin1980} and AKS \citep{agrawal2004} are
standard primality-testing references. This work does not study primality
testing itself; it studies candidate-generation policies.

\paragraph{Learned heuristics for mathematical search.}
Machine-learning approaches to mathematical discovery and heuristic search
have been explored in recent work such as \citet{davies2021}. Our learned
bit-order policy is much simpler: it is a count-based heuristic rather than a
neural model.

%% ---------------------------------------------------------------
\section{Conclusion}
\label{sec:conclusion}
%% ---------------------------------------------------------------

We studied prime numbers as vertices in the binary hypercube and compared
Hamming-space prime search with standard arithmetic prime search.

The key empirical findings, for the reported experimental settings, are:

\begin{enumerate}
  \item \textbf{Hamming density.} In the tested small dimensions, nearly every
    integer is within a small number of bit flips of a prime, consistent with
    the rapid growth of Hamming balls.

  \item \textbf{Parity dominates the full hypercube.} In unrestricted Hamming
    search, bit~0 is overwhelmingly useful because it changes parity. This
    motivates the odd-subcube correction, where bit~0 is fixed and only odd
    candidates are considered.

  \item \textbf{SOI wins average checks; LBF is competitive per-instance.}
    Standard odd-increment search achieves the lowest average candidate checks
    in all tested configurations. However, low-bit-first odd-subcube search
    wins or ties SOI on roughly 58--60\% of individual inputs at
    $m \in \{24,28,32\}$.

  \item \textbf{Two geometries.} Standard search optimizes arithmetic locality;
    Hamming search optimizes bit-space locality. Pareto analysis shows that
    neither notion of locality fully subsumes the other.
\end{enumerate}

\subsection{Limitations}

This work is empirical. The reported experiments cover moderate bit sizes and
a finite number of random samples. We do not prove asymptotic claims about
nearest-prime Hamming distance, win rates, or Pareto structure. In addition,
the Hamming strategies tested here are simple orderings of bit-flip masks;
more sophisticated modular-residue-aware or learned policies may change the
observed tradeoffs. Finally, the work does not claim to improve standard
next-prime search for arithmetic-local prime generation.

\subsection{Open questions}

\begin{enumerate}[label=(\roman*)]
  \item \textbf{Larger bit sizes.} Do the win rates and Pareto structure hold
    at $m \in \{40,48,64\}$ and beyond?

  \item \textbf{Multiple seeds and confidence intervals.} How stable are the
    reported win rates and frontier positions across seeds?

  \item \textbf{Hamming-1 prime pairs.} How many pairs of primes $p,q$ satisfy
    $\dH(p,q)=1$, i.e., differ by one power of two? What is their distribution
    across bit positions?

  \item \textbf{Hypercube sieve.} Can modular structure be used to prescreen
    Hamming candidates, analogous to a sieve in arithmetic space?

  \item \textbf{Modular-residue-aware scoring.} Can bit positions be scored by
    their direct effect on residues modulo small primes, rather than learned
    from data?

  \item \textbf{Improved hybrid policies.} Can arithmetic moves and Hamming
    moves be combined into a better multi-objective search policy?
\end{enumerate}

\subsection*{Code availability}

All code, data generation scripts, analysis, and blog chapter drafts are
available at:
\[
  \text{\url{https://github.com/pojhafb/prime-hypercube-search}}
\]

%% ---------------------------------------------------------------
\appendix
%% ---------------------------------------------------------------

\section{Flip-set construction}
\label{app:flipsets}

For $m$ bits and maximum radius $r_{\max}$, the full list of flip sets in the
odd subcube is:
\[
  \mathcal{F}(m,r_{\max})
  =
  \bigcup_{r=0}^{r_{\max}}
  \binom{\{1,\ldots,m-1\}}{r},
\]
where $\binom{S}{r}$ denotes all $r$-element subsets of $S$.

The total number of flip sets is $B(m-1,r_{\max})$.

\begin{center}
\begin{tabular}{rrrrrr}
\toprule
$m$ & $r=0$ & $r=1$ & $r=2$ & $r=3$ & $r=4$ total \\
\midrule
16 & 1 & 15 & 105 & 455 & 1941 \\
32 & 1 & 31 & 465 & 4495 & 36457 \\
64 & 1 & 63 & 1953 & 39711 & 637393 \\
\bottomrule
\end{tabular}
\end{center}

At $m=64$, the radius-4 odd-subcube flip-set list contains 637,393 candidates.
Caching per $(m,\text{strategy})$ reduces this to a one-time cost per experiment.

\section{Primality testing}
\label{app:primality}

We use \texttt{sympy.isprime}, which combines trial division and primality tests
such as Miller--Rabin and Baillie--PSW depending on the size of the input.
Results are cached via \texttt{functools.lru\_cache} with capacity
$2\times 10^6$.

The goal of this work is not to benchmark primality-testing algorithms
themselves, but to compare candidate-generation policies. All strategies in a
given experiment use the same primality-testing backend.

%% ---------------------------------------------------------------
\begin{thebibliography}{99}

\bibitem[Agrawal et~al.(2004)]{agrawal2004}
Agrawal, M., Kayal, N., and Saxena, N. (2004).
\newblock PRIMES is in P.
\newblock \emph{Annals of Mathematics}, 160(2), 781--793.

\bibitem[Cram\'er(1936)]{cramer1936}
Cram\'er, H. (1936).
\newblock On the order of magnitude of the difference between consecutive prime numbers.
\newblock \emph{Acta Arithmetica}, 2, 23--46.

\bibitem[Davies et~al.(2021)]{davies2021}
Davies, A., Veli\v{c}kovi\'c, P., Buesing, L., et~al. (2021).
\newblock Advancing mathematics by guiding human intuition with AI.
\newblock \emph{Nature}, 600, 70--74.

\bibitem[de~la Vall\'ee~Poussin(1896)]{delavalleepoussin1896}
de~la Vall\'ee~Poussin, C.~J. (1896).
\newblock Recherches analytiques sur la th\'eorie des nombres premiers.
\newblock \emph{Annales de la Soci\'et\'e Scientifique de Bruxelles}, 20, 183--256.

\bibitem[Granville(1995)]{granville1995}
Granville, A. (1995).
\newblock Harald Cram\'er and the distribution of prime numbers.
\newblock \emph{Scandinavian Actuarial Journal}, 1995(1), 12--28.

\bibitem[Hadamard(1896)]{hadamard1896}
Hadamard, J. (1896).
\newblock Sur la distribution des z\'eros de la fonction $\zeta(s)$ et ses
  cons\'equences arithm\'etiques.
\newblock \emph{Bulletin de la Soci\'et\'e Math\'ematique de France}, 24, 199--220.

\bibitem[Knuth(2011)]{knuth2011}
Knuth, D.~E. (2011).
\newblock \emph{The Art of Computer Programming, Volume~4A: Combinatorial
  Algorithms, Part~1}.
\newblock Addison-Wesley.

\bibitem[MacWilliams and Sloane(1977)]{macwilliams1977}
MacWilliams, F.~J. and Sloane, N.~J.~A. (1977).
\newblock \emph{The Theory of Error-Correcting Codes}.
\newblock North-Holland.

\bibitem[Maynard(2015)]{maynard2015}
Maynard, J. (2015).
\newblock Small gaps between primes.
\newblock \emph{Annals of Mathematics}, 181(1), 383--413.

\bibitem[Menezes et~al.(1996)]{menezes1996}
Menezes, A.~J., van Oorschot, P.~C., and Vanstone, S.~A. (1996).
\newblock \emph{Handbook of Applied Cryptography}.
\newblock CRC Press.

\bibitem[Miller(1976)]{miller1976}
Miller, G.~L. (1976).
\newblock Riemann's hypothesis and tests for primality.
\newblock \emph{Journal of Computer and System Sciences}, 13(3), 300--317.

\bibitem[O'Donnell(2014)]{odonnell2014}
O'Donnell, R. (2014).
\newblock \emph{Analysis of Boolean Functions}.
\newblock Cambridge University Press.

\bibitem[Rabin(1980)]{rabin1980}
Rabin, M.~O. (1980).
\newblock Probabilistic algorithm for testing primality.
\newblock \emph{Journal of Number Theory}, 12(1), 128--138.

\end{thebibliography}

%% ---------------------------------------------------------------
\end{document}
